\documentclass[11pt]{amsart}
\usepackage[localtheorems]{raeez-math-template}

\newtheorem{theorem}{Theorem}[section]
\newtheorem{proposition}[theorem]{Proposition}
\newtheorem{lemma}[theorem]{Lemma}
\newtheorem{corollary}[theorem]{Corollary}
\theoremstyle{definition}
\newtheorem{definition}[theorem]{Definition}
\newtheorem{example}[theorem]{Example}
\theoremstyle{remark}
\newtheorem{remark}[theorem]{Remark}

\newcommand{\En}{\mathsf{E}_n}
\newcommand{\Eone}{\mathsf{E}_1}
\newcommand{\Etwo}{\mathsf{E}_2}
\newcommand{\Ethree}{\mathsf{E}_3}
\newcommand{\Einf}{\mathsf{E}_\infty}
\newcommand{\Ezero}{\mathsf{E}_0}
\newcommand{\Rep}{\mathrm{Rep}}
\newcommand{\HH}{\mathrm{HH}}
\providecommand{\RHom}{\mathrm{RHom}}
\providecommand{\PhiFA}{\Phi^{\mathrm{FA}}}
\providecommand{\SpCh}{\mathrm{Sp}^{\mathrm{ch}}}
\providecommand{\Obs}{\mathrm{Obs}}
\newcommand{\cZ}{\mathcal{Z}}
\newcommand{\cC}{\mathcal{C}}
\newcommand{\C}{\mathbb{C}}
\newcommand{\Z}{\mathbb{Z}}
\newcommand{\fgl}{\mathfrak{gl}}
\newcommand{\PV}{\mathrm{PV}}
\newcommand{\CoHA}{\mathrm{CoHA}}
\newcommand{\cW}{\mathcal{W}}

\title{The $\En$ Hierarchy, the $\Eone$ Mechanism, and the Centre}
\author{Raeez Lorgat}
\date{}

\begin{document}
\maketitle

\begin{abstract}
The dimensional stratification of the CY-to-chiral correspondence has
three independent inputs.  Dunn-Lurie additivity is an equivalence of
$\infty$-operads; restriction along a little-line inclusion gives an
$\Etwo$-to-$\Eone$ forgetful functor.  The $\Omega$-background is the
toric equivariant mechanism for local CY$_3$ geometries.  Compact
non-toric CY$_3$ inputs are governed by the framed specialisation
theorem and its chain-level hypotheses.  On constructed $d\geq 3$ loci,
the specialised algebra $A$ is $\Eone$-chiral.  The $\Etwo$ braiding
lives on a centre, double, or representation-theoretic enhancement once
the pairing and completion data exist.  The Hall, shuffle, factorisation,
Chern-Simons, and automorphic constructions around $\C^3$ and
$K3\times E$ have different inputs; comparison theorems relate their
outputs.
\end{abstract}

\section{Conventions}

All operads are little-disks operads in spectra or in chain complexes
over $\C$.  Dunn-Lurie additivity is the equivalence
\begin{equation}
  \Etwo \simeq \Eone\otimes_{\Ezero}\Eone
  \label{eq:dunn}
\end{equation}
in the Boardman-Vogt tensor product of $\infty$-operads, and more
generally $\En\simeq\Eone^{\otimes n}$.  Thus an $\Etwo$-algebra is an
$\Eone$-algebra internal to $\Eone$-algebras:
\begin{equation}
  \mathsf{Alg}_{\Etwo}(\cC)
  \simeq
  \mathsf{Alg}_{\Eone}\bigl(\mathsf{Alg}_{\Eone}(\cC)\bigr).
  \label{eq:dunn-iterated}
\end{equation}

The output convention is:
\[
\begin{array}{c|c|c}
d & \Phi_d\text{-output scope} & \text{basic source} \\
\hline
1 & E_\infty\text{-chiral} & \text{elliptic and lattice VOA branch}\\
2 & E_2\text{-chiral} & \text{K3 Mukai-Heisenberg branch}\\
\geq 3 & E_1\text{-chiral on the curve} &
\text{framed specialisation and local Hall branches}
\end{array}
\]
For $d\geq 3$, the $E_2$-braided object is
$\cZ(\Rep^{E_1}(A))$ or a Drinfeld double on constructed loci.  The
native algebra $A$ remains $E_1$.

The notation
\[
  \Phi_d^{(\Sigma_{d-1},C)}
  =
  \SpCh_{\Sigma_{d-1},C}\circ\PhiFA_d
\]
has object-level meaning only after the formality, framing, anomaly,
completion, and witnessed specialisation data have been fixed.  Its
morphism-level functoriality is restricted to kernels preserving those
data.  The family $\{\Phi_d\}_{d\geq1}$ has dimension-dependent targets;
the superscript \((\Sigma_{d-1},C)\) is part of the datum at $d\geq3$.

\section{The Five Objects Attached to an \texorpdfstring{$E_1$}{E1} Algebra}
\label{sec:five-objects}

Let $A$ be an $E_1$-chiral algebra on a curve.  Five functorial objects
occur in the bar-cobar and centre formalism:
\[
  A,\qquad
  B(A),\qquad
  A^i,\qquad
  A^!,\qquad
  Z^{\mathrm{der}}_{\mathrm{ch}}(A).
\]
Here $B(A)$ is the chiral bar coalgebra.  The equation
\[
  \Omega B(A)\simeq A
\]
is bar-cobar inversion under the connectivity and completion
hypotheses.  The object $A^i$ is the Koszul-dual coalgebra supplied by
the cooperadic bar construction.  The chiral algebra $A^!$ is the
Verdier/Koszul dual, obtained from the Verdier dual of the bar coalgebra
when the finiteness hypotheses needed for linear duality hold.  The
bulk object is
\[
  Z^{\mathrm{der}}_{\mathrm{ch}}(A)
  =
  \RHom^{\mathrm{ch}}_{A\text{-bimod}}(A,A)
  \simeq
  C^\bullet_{\mathrm{ch}}(A,A),
\]
the chiral Hochschild cochain algebra.  Bar-cobar inversion, Koszul
duality, Verdier duality, and Hochschild centre formation are four
different constructions.

\section{Dunn Additivity and Forgetful Directions}

\begin{proposition}[Operadic additivity]
\label{prop:dunn-scope}
The equivalence \eqref{eq:dunn} is an equivalence of
$\infty$-operads.  A choice of embedding
$\Eone\to\Etwo$ along one coordinate direction gives a pullback functor
\[
  \mathsf{Alg}_{\Etwo}(\cC)
  \longrightarrow
  \mathsf{Alg}_{\Eone}(\cC).
\]
This functor remembers the selected multiplication and forgets the
coherent transverse multiplication.
\end{proposition}

\begin{proof}
Dunn-Lurie additivity identifies $\Etwo$-algebras with
$\Eone$-algebras in $\Eone$-algebras.  Restricting an
$\Etwo$-algebra along one little-line inclusion of the little-disk
operad keeps the multiplication in that direction and forgets the
coherent transverse multiplication.  The construction is pullback along
a map of operads.  A tensor-factor projection would require a cartesian
product presentation, while \eqref{eq:dunn} uses the Boardman-Vogt
tensor product.
\end{proof}

\section{The $\Omega$-Background in Toric and Compact Cases}
\label{sec:omega}

\begin{proposition}[Toric local mechanism]
\label{prop:omega-toric}
Let $X$ be a toric local CY$_3$ carrying a two-torus action on the
transverse complex plane with equivariant parameters
$\epsilon_1,\epsilon_2$ and Calabi-Yau constraint
$\epsilon_1+\epsilon_2+\epsilon_3=0$.  Equivariant localisation of the
factorisation envelope along the transverse directions produces an
$E_1$-chiral algebra on the remaining curve direction.  This is the
Costello-Li mechanism for the local $\C^3$ and related toric examples.
\end{proposition}

\begin{proof}
The $\Omega$-background is an equivariant deformation by the chosen
torus action.  Localisation removes the moving transverse modes and keeps
the fixed curve direction.  Operadically, this is the pullback described
in Proposition~\ref{prop:dunn-scope}, selected by the geometric fixed
locus inside the toric geometry.
\end{proof}

\begin{proposition}[Compact CY$_3$ mechanism]
\label{prop:compact-cy3-mechanism}
For a compact non-toric CY$_3$, the toric equivariant deformation is
absent.  The available object-level construction is the framed
specialisation
\[
  \Phi_3^{(\Sigma_2,C)}(\cC)
  =
  \SpCh_{\Sigma_2,C}\bigl(\PhiFA_3(\cC)\bigr)
\]
under the hypotheses of the CY$_3$ specialisation theorem: smoothness,
properness, a negative-cyclic Calabi-Yau class with non-degenerate
Hochschild shadow, a chain-level $\mathbb{S}^3$-framing compatible with
the BV operator, the TCFT correction, anomaly cancellation, analytic
completion, and the filtered $\HH^{-2}_{E_1}(A,A)$ obstruction
comparison.  Its native output is $E_1$-chiral on $C$.
\end{proposition}

\begin{proof}
The construction uses the witnessed specialisation datum and the
chain-level hypotheses listed in the statement.  The stage-$1$
Hochschild calculus
has Gerstenhaber bracket degree $1-d$ and PTVV shift $2-d$.  The
stage-$2$ restriction along $(\Sigma_2,C)$ leaves a curve-direction
chiral algebra.  At $d=3$ this residual object is $E_1$; the $E_2$
braiding is obtained from the centre of the representation category on
the loci where the centre, pairing, and completion are constructed.
\end{proof}

\section{Distinct Constructions Around \texorpdfstring{$\C^3$}{C3}}
\label{sec:routes-c3}

Five standard constructions enter the local $\C^3$ comparison, each
with its own input.
\begin{enumerate}[label=\textup{(R\arabic*)}]
\item \emph{Critical CoHA.}  The $3$-Calabi-Yau quiver or
matrix-factorisation input produces
$\CoHA(\C^3)\simeq Y^+(\widehat{\fgl}_1)$ by the
Schiffmann-Vasserot comparison.  This is an $E_1$ positive half.
\item \emph{Plane partitions.}  The MacMahon Fock space carries the
level-one affine-Yangian action.  The $\cW_{1+\infty}$ vacuum module is
obtained from the full Yangian action after the positive half has been
completed.
\item \emph{Shuffle algebra.}  Negut's shuffle presentation gives the
same positive half $Y^+(\widehat{\fgl}_1)$ with its explicit structure
function.
\item \emph{Factorisation envelope.}  The dg Lie algebra
$\PV^\ast(\C^3)[2]$ with Schouten bracket has a
Costello-Gwilliam chiral envelope; after the toric specialisation it
matches the same local positive-half branch.
\item \emph{Holomorphic-topological Chern-Simons.}  The perturbative
boundary algebra of the Costello-Li local theory is compared with the
completed double and its $\cW_{1+\infty}$ vacuum representation after
the Hopf pairing, completion, and evaluation-module data are supplied.
\end{enumerate}

The Hall and shuffle comparisons identify the positive half.  The direct
Hall output on $\C^3$ is
\[
  \CoHA(\C^3)=Y^+(\widehat{\fgl}_1).
\]
The doubled and evaluated layer is typed by the chain
\[
  Y^+(\widehat{\fgl}_1)
  \longrightarrow
  D^\wedge\!\bigl(Y^+(\widehat{\fgl}_1)\bigr)
  =
  Y(\widehat{\fgl}_1)
  \longrightarrow
  \mathrm{End}\bigl(\cW_{1+\infty}\text{-}\mathrm{vac}\bigr),
\]
with the first arrow requiring the Hopf pairing and completion.

\begin{remark}[K3 \texorpdfstring{$\times$}{x} E]
The analogous $K3\times E$ constructions are relative DT theory,
Borcherds-Gritsenko BKM, Maulik-Okounkov stable envelopes,
Vafa-Witten theory, and the relative chiral algebra.  They have
different inputs and comparison maps.  The numerical package consists
of five measurements with four distinct values:
\[
\begin{aligned}
 \kappa_{\mathrm{cat}}(K3\times E)
   &=\chi(\mathcal O_{K3})\,\chi(\mathcal O_E)=2\cdot0=0,\\
 \kappa_{\mathrm{ch}}^{\mathrm{Hodge}}(K3\times E)
   &=\sum_q(-1)^q h^{0,q}(K3\times E)=0,\\
 \kappa_{\mathrm{ch}}^{\mathrm{Heis}}(K3\times E)&=3,\\
 \kappa_{\mathrm{BKM}}(\Delta_5)&=c_1(0)/2=5,\\
 \kappa_{\mathrm{fiber}}(K3)&=24.
\end{aligned}
\]
\end{remark}

\section{The Centre is the Source of the \texorpdfstring{$E_2$}{E2} Braiding}
\label{sec:centre}

\begin{proposition}[Derived centre]
\label{prop:centre}
For an $E_1$-algebra $A$ in a presentable symmetric monoidal stable
category with internal bimodule Homs,
\[
  Z^{\mathrm{der}}_{E_1}(A)
  =
  \mathrm{End}^\bullet_{A\text{-bimod}}(A)
  \simeq
  \HH^\bullet(A,A)
\]
is an $E_2$-algebra.  In the chiral setting this is the underlying
topological form of
\[
  Z^{\mathrm{der}}_{\mathrm{ch}}(A)
  =
  C^\bullet_{\mathrm{ch}}(A,A).
\]
Categorically this centre controls the Drinfeld centre of
$\Rep^{E_1}(A)$.
\end{proposition}

\begin{proof}
The Deligne conjecture gives an $E_2$ action on Hochschild cochains.
Equivalently, the centre is the endomorphism object of the identity
bimodule, and the two-dimensional little-disk operations act by
composition and interchange of bimodule endomorphisms.  At the
categorical level this is the half-braiding definition of the
Drinfeld centre.
\end{proof}

\begin{remark}[Right-adjoint character]
The centre is a right-adjoint type construction: it is an endomorphism
object and commutes with limits under the usual finiteness hypotheses.
Averaging over a group action produces invariants or coinvariants of the
original $E_1$ multiplication.  Centre formation takes endomorphisms of
the identity bimodule; the interchange of these endomorphisms supplies
the second multiplication whose hexagon equations give an $E_2$
braiding.
\end{remark}

\begin{remark}[Local \texorpdfstring{$\C^3$}{C3}]
For the local Hall branch,
\[
  \cZ\bigl(\Rep^{E_1}(Y^+(\widehat{\fgl}_1))\bigr)
\]
is the representation-categorical place where the full affine Yangian
braiding appears after the Hopf pairing and completion are fixed.  The
$R$-matrix is a chain-level family
$R(u)\in\mathrm{End}(V(u)\otimes V(v))[[u^{-1}]]$ satisfying the
hexagon/Yang-Baxter identities.  The spectral parameter comes from
evaluation modules for the loop algebra; the centre supplies the
half-braiding equations.
\end{remark}

\section{Framing Data and Operadic Level}
\label{sec:framing}

Framing obstructions and operadic enrichment are orthogonal.  A higher
$E_n$ structure is extra multiplication data.  A framing obstruction is
data needed to integrate or specialise a fixed factorisation object.

There are two Bott towers.  The complex tower is
$\pi_d(BU)=\Z$ for even $d\geq2$ and $0$ for odd $d$.  The symplectic
refinement is used on odd CY dimension, where the Serre pairing is
antisymmetric.  In the range relevant here the effective obstruction is
\[
\begin{array}{c|cccccc}
d & 3 & 4 & 5 & 6 & 7 & 8\\
\hline
\pi_d(BU) & 0 & \Z & 0 & \Z & 0 & \Z\\
\pi_d(B\mathrm{Sp})\text{ on odd }d & 0 & - & \Z/2 & - & 0 & -\\
\Obs_{\mathrm{eff}}(d) & 0 & \Z & \Z/2 & \Z & 0 & \Z
\end{array}
\]
The $BU$ column has period \(2\), while the real and symplectic
refinements have period \(8\).  These groups record grading, torsion, or
twist data in the framed specialisation.  The native operadic level of
the specialised output remains $E_1$ for $d\geq3$.

\begin{proposition}[CY$_3$ framing data]
\label{prop:cy3-framing-obligations}
The vanishing of the stable topological class at $d=3$ is only the
first obstruction check.  A compact CY$_3$ object-level construction
still requires the chain-level $\mathbb{S}^3$-framing homotopy, TCFT
correction, Costello-Li anomaly cancellation, analytic completion, and
the filtered $\HH^{-2}_{E_1}(A,A)$ comparison.  Once those data are
present, the specialised output is $E_1$ and its centre carries the
$E_2$ braiding.
\end{proposition}

\begin{proof}
The topological obstruction class controls the existence of a framing
at the stable topological level.  The construction is chain-level and
quantum.  The TCFT correction and Costello-Li completion control the
effective BV differential; the filtered
Hochschild comparison removes the primary derived obstruction named by
the CY$_3$ specialisation theorem.  Higher lifting-space layers remain
outside this statement.  These steps supply the data required for the
$E_1$ specialisation; the second multiplication belongs to the centre construction of
Proposition~\ref{prop:centre}.
\end{proof}

\section{Factorisation-Homological Form of the Statement}
\label{sec:reconstruction}

\begin{theorem}[Specialised CY$_3$ output]
\label{thm:e1-output-d3}
Assume the hypotheses of the CY$_3$ specialisation theorem for a CY$_3$
category $\cC$ and a chosen pair $(\Sigma_2,C)$.  Then
\[
  A_\cC
  =
  \Phi_3^{(\Sigma_2,C)}(\cC)
  =
  \SpCh_{\Sigma_2,C}\bigl(\PhiFA_3(\cC)\bigr)
\]
is an $E_1$-chiral algebra on the curve $C$.  If the representation
category, Hopf pairing, centre, and completion data are constructed,
$\cZ(\Rep^{E_1}(A_\cC))$ is $E_2$-braided.
\end{theorem}

\begin{proof}
The first stage $\PhiFA_3$ is controlled by the
degree-$-2$ Gerstenhaber bracket on Hochschild cochains and the
corresponding shifted symplectic form.  The specialisation
$\SpCh_{\Sigma_2,C}$ restricts the factorisation object along
the chosen surface and leaves a curve-direction chiral algebra.  A
curve chiral algebra is $E_1$ in the present convention.  The centre
statement is Proposition~\ref{prop:centre}.
\end{proof}

\begin{remark}[Higher dimensions]
For $d\geq 4$ the same specialisation lane is conditional outside
constructed loci.  Bott and framing classes specify grading and twisting
data for the prospective $E_1$ output.  Native higher multiplication
requires an additional centre, double, module, or enhancement
construction.
\end{remark}

\section{Hypotheses Beyond Constructed Loci}

The general compact and higher-dimensional statements require the
following inputs.
\begin{enumerate}[label=\textup{(\arabic*)}]
\item For an arbitrary compact non-formal CY$_3$ category, construct a
strictified Hochschild $\Eone$ model, a compatible negative cyclic
Calabi-Yau class, the chain-level $\mathbb{S}^3$-framing homotopy, the
TCFT correction, anomaly cancellation, and the holomorphic completion
needed for $\SpCh_{\Sigma_2,C}$.
\item For the centre comparison, prove the smoothness, properness, and
dualizability hypotheses that identify
$\cZ(\Rep^{\Eone}(A))$ with
$\Rep^{\Etwo}(Z^{\mathrm{der}}_{\mathrm{ch}}(A))$.
\item For a Drinfeld double, supply the non-degenerate Hopf or
stable-envelope pairing, the topological completion, and compatibility
of the completed double with the chiral centre.
\item For $K3\times E$, compare the relative DT, Borcherds-Gritsenko,
Maulik-Okounkov, Vafa-Witten, and relative chiral constructions by
explicit maps of their input categories or representation categories;
character equality alone is insufficient.
\item For $d\geq4$, construct the analogue of the framed
specialisation theorem before asserting a compact Stage-$2$ output.
Higher operadic structure belongs to a centre, double, module, or
enhancement layer.
\end{enumerate}

\section{Operadic Consequences}

\begin{corollary}
The CY-to-chiral hierarchy obeys the following dimension and enhancement
rules:
\[
\begin{array}{lll}
\text{input} & \text{theorem or construction} & \text{conclusion}\\
\hline
\text{Dunn additivity}
& \En\simeq\Eone^{\otimes n}
& \text{restriction gives forgetful } E_n\to E_1\\
\Omega\text{-background}
& \text{toric equivariant localisation}
& \text{local CY}_3\text{ route to }E_1\\
\text{compact CY}_3
& \Phi_3^{(\Sigma_2,C)}=\SpCh(\PhiFA_3)
& E_1\text{ under the framed hypotheses}\\
\text{centre}
& Z^{\mathrm{der}}_{E_1}(A)=\HH^\bullet(A,A)
& E_2\text{ braiding on }\cZ(\Rep^{E_1}(A))\\
\text{framing}
& \Obs_{\mathrm{eff}}(d)\text{ from }BU/B\mathrm{Sp}
& \text{grading or twist of the }E_1\text{ object}
\end{array}
\]
\end{corollary}

\end{document}
